% !TSS name: Стаття з standalone-рисунком (папка: .tex + fig1.tex)
% !TSS description: шаблон із двох файлів: стаття (LuaLaTeX, кирилиця) підключає fig1.pdf, а fig1.tex це standalone TikZ-рисунок; спершу збери fig1.tex окремо (latexmk -lualatex fig1.tex), потім статтю (питає автора)
% !TeX program = lualatex
% !TeX encoding = utf8
% !TeX spellcheck = uk_UA
\documentclass[a4paper, 12pt]{article}
\usepackage{fontsetup}
\usepackage[english, ukrainian]{babel}
\usepackage[a4paper, margin=2cm]{geometry}
\usepackage{graphicx}
\usepackage{amsmath}
\usepackage{hyperref}

\title{}
\author{${ask:author}}

\begin{document}
\maketitle

Рисунок~\ref{fig:one} зроблено окремим документом \texttt{fig1.tex} (клас \texttt{standalone}),
який збирається у \texttt{fig1.pdf}.

\begin{figure}[h!]
	\centering
	\includegraphics[width=0.6\linewidth]{fig1.pdf}
	\caption{${cursor}}\label{fig:one}
\end{figure}

\end{document}
